% Copyright 2026 Yoshida Shin
%
% This file is part of the ``Repeating Nines''.
%
% Permission is granted to copy, distribute, and/or modify this document
% under the terms of the GNU Free Documentation License, Version 1.3 or any later version published by the Free Software Foundation;
% with no Front-Cover Texts, no Back-Cover Texts, and one Invariant Section: ``This document was written by Shin Yoshida with the aim of contributing to a fairer and safer world.''
% A copy of the license is included in the section entitled ``GNU Free Documentation License.''


\begin{frame}[c]{}{}
  {\large Intuitive Understanding}
\end{frame}

\begin{frame}[c]{}{}
  Let's consider the following number sequence:

  \vspace{2ex}

  $a_n = 1 - (0.1)^n$

  \vspace{2ex}

  i.e.

  $a_1$ = 0.9

  $a_2$ = 0.99

  $a_3$ = 0.999

  ...
\end{frame}

\begin{frame}[c]{}{}
  Do you see that $a_n$ has the following characteristics?

  \vspace{2ex}

  \begin{itemize}
  \item $\forall n \in \mathbb{N}, a_n < a_{n+1}$
  \item $\forall n \in \mathbb{N}, a_n < 10$
  \end{itemize}
\end{frame}

\begin{frame}[c]{}{}
  You might think,

  ``The second one is wrong.

  It should be $\forall n \in \mathbb{N}, a_n < 1$.''

  \vspace{2ex}

  That is correct, but at the same time, ``$a_n < 10$'' is also correct.
\end{frame}

\begin{frame}[c]{}{}
  $\forall n \in \mathbb{N}, a_n < 10$

  \vspace{4ex}

  This is wrong only if

  \vspace{2ex}

  $\exists n \in \mathbb{N}, 10 \leqq a_n$
\end{frame}

\begin{frame}[c]{}{}
  I don't think these are so obvious that we need to prove them, but this isn't a math test.

  \vspace{2ex}

  Let's skip the proof.
\end{frame}

\begin{frame}[c]{}{}
  Next, let's think about what the graph of $a_n$ looks like.

  \vspace{4ex}

  Btw, my explanation in this section is not mathematically rigorous. Try to find the flaw!
\end{frame}

\begin{frame}[c]{}{}
  As I mentioned previously,

  \vspace{2ex}

  $a_n$ fulfills the following condition.

  \vspace{4ex}

  $\forall n \in \mathbb{N}, a_n < a_{n+1}$ 
\end{frame}

\begin{frame}[c]{}{}
 Then, the graph of $a_n$ must clearly look like either of the following 2 options.
\end{frame}

\begin{frame}[c]{}{}
  Option 1

  \vspace{2ex}

  $a_n$ increases without an upper limit as n gets larger.

  \begin{center}
    \includegraphics[width=0.5\textwidth]{graphs/graph0.png}
  \end{center}
\end{frame}

\begin{frame}[c]{}{}
  Option 2

  \vspace{2ex}

  $a_n$ increases as n gets larger, but there is an upper limit.

  \begin{center}
    \includegraphics[width=0.5\textwidth]{graphs/graph1.png}
  \end{center}
\end{frame}

\begin{frame}[c]{}{}
  Now, remember the other condition:

  \vspace{2ex}

  $\forall n \in \mathbb{N}, a_n < 10$

  \vspace{4ex}

  This obviously conflicts with option 1.
\end{frame}

\begin{frame}[c]{}{}
  So, the graph of $a_n$ must be like option 2:

  \vspace{4ex}

  $a_n$ increases as n gets larger, but there is an upper limit.
\end{frame}

\begin{frame}[c]{}{}
  Btw, did you find the flaw?

  \vspace{4ex}

  The answer is in the next page.
\end{frame}

\begin{frame}[c]{}{}
  I should have defined the term ``upper limit.''

  \vspace{4ex}

  But this section is just for intuitive understanding.

  \vspace{2ex}

  Let me skip that.
\end{frame}

\begin{frame}[c]{}{}
  Intuitively, you can understand that

  \vspace{2ex}

  $a_n$ gets closer and closer to some value as n gets larger.

  \vspace{2ex}

  I want to call this value ``0.9999...''

  \begin{center}
    \includegraphics[width=0.5\textwidth]{graphs/graph2.png}
  \end{center}
\end{frame}

\begin{frame}[c]{}{}
  In other words, ``0.9999...'' should be the value that is just slightly larger than any number ``0.'' followed by a lot of ``9s''.
\end{frame}

\begin{frame}[c]{}{}
  This value is called ``Limit'' in mathematics.

  \vspace{2ex}

  We usually write it as $$\lim_{n \to \infty} a_n$$.
\end{frame}
